\section{Shear--locking--free FE--Meshfree formulation}
\subsection{Reproducing kernel meshfree approximation}
\begin{equation}
    \Psi_I(\boldsymbol x) = \boldsymbol c^T(\boldsymbol x_I-\boldsymbol x)\boldsymbol p(\boldsymbol x_I-\boldsymbol x)\phi(\boldsymbol x_I-\boldsymbol x)
\end{equation}
\begin{equation}
    \phi(\boldsymbol x_I-\boldsymbol x) = \phi(r_{1I})\phi(r_{2I}),\quad
    r_{iI} = \frac{\vert x_{iI}-x_i\vert}{s_i}
\end{equation}
\begin{equation}
    \phi(r) = \frac{1}{3!}\left\{
        \begin{aligned}
            &(2-2r)^3-4(1-2r)^3 \quad& r\le \frac{1}{2} \\
            &(2-2r)^3 & \frac{1}{2} < r \le 1 \\
            &0 & r > 1
        \end{aligned}
    \right .
\end{equation}
\begin{equation}
    \sum_{I=1}^{n} \Psi_I(\boldsymbol x)\boldsymbol p(\boldsymbol x_I) = \boldsymbol p(\boldsymbol x)
\end{equation}
or a shifted form
\begin{equation}
    \sum_{I=1}^n \Psi_I(\boldsymbol x) \boldsymbol p(\boldsymbol x_I-\boldsymbol x) = \boldsymbol p(\boldsymbol 0)
\end{equation}
\begin{equation}
    \boldsymbol c(\boldsymbol x_I-\boldsymbol x) = \boldsymbol A^{-1}(\boldsymbol x_I-\boldsymbol x)\boldsymbol p(\boldsymbol 0)
\end{equation}
with
\begin{equation}
    \boldsymbol A(\boldsymbol x_I-\boldsymbol x) = \sum_{I=1}^n \boldsymbol p(\boldsymbol x_I-\boldsymbol x)\boldsymbol p^T(\boldsymbol x_I-\boldsymbol x)\phi(\boldsymbol x_I-\boldsymbol x)
\end{equation}
\begin{equation}
    \Psi_I(\boldsymbol x) = \boldsymbol p^T(\boldsymbol 0) \boldsymbol A^{-1} \boldsymbol p(\boldsymbol x_I-\boldsymbol x)\phi(\boldsymbol x_I-\boldsymbol x)
\end{equation}
\subsection{Eigenvalue test for checking LBB coefficients(should be fixed)}
To verify the optimal constraint ratios proposed in Eq. \eqref{eq:opt_n},
the eigenvalue test of LBB coefficients \cite{bathe2001} in a square domain, shown in Figure.\ref{}, is introduced herein, where all the boundaries of this domain are clamped.
In order to adjust the constraint ratio more flexibly, 
the reproducing kernel meshfree approximation are used to discretize the deflection $w$ and rotation $\boldsymbol \varphi$, yields:
\begin{equation}
    N^w_I := \Psi_I,\; N^\varphi_K := \Psi_K^\varphi,\; N^q_R := N_R
\end{equation}
while, as shown in Figure.\ref{}, the meshfree nodes for $w$ and $\boldsymbol \varphi$ are different. 

To meet the variational consistency of the weak formulation with meshfree approximation, the reproducing kernel gradient smoothing integration framework (RKGSI) \cite{wu2023} is used herein. 
Accordingly, the Mindlin plate version of RKGSI framework is presented in \ref{sec:rkgsi}.
\subsection{Shear--locking--free formulations with optimal constraint ratio}

\begin{equation}
    N^w_I := \Psi_I,\; N^\varphi_K := N_K,\; N^q_R := N_R
\end{equation}
FIXME
\begin{rmk}
The discretization of Eq. 
\end{rmk}